\section{The translation}
\label{sec:translation}

This section defines the target theory.  Throughout, fix a well-formed
definitional environment $\Env$ whose declarations lie in $\Sfrag$
(\cref{def:shallow-fragment}) and a closed shallow goal proposition $G$.

The target logic is classical first-order logic with equality~$\feq$; we
write $\provesfol$ for provability in any standard sound and complete
proof calculus \cite{Shoenfield1967}.  We assume the propositional
constants $\top$ and $\bot$ are available in the target language (they can
be read as abbreviations).

\subsection{Signature}
\label{sec:translation-sig}

\begin{definition}[Signature $\SigFOL$]
\label{def:signature}
The signature $\SigFOL$ contains:
\begin{itemize}[leftmargin=2em]
\item a unary predicate symbol $\Ihat{I}$ (the \emph{guard predicate}) for
  each datatype $I$ of $\Env$;
\item a $k$-ary predicate symbol $\Ihat{J}$ for each inductive predicate
  $J : (\sigma_1, \dots, \sigma_k) \to \Prop$ of $\Env$;
\item an $r$-ary function symbol $\Ihat{\wcns{c}}$ for each constructor
  $\wcns{c} : (\sigma_1, \dots, \sigma_r)$ of a datatype of $\Env$;
\item a function symbol $\Ihat{c}$ of arity $a_c$, and if $a_c > 0$
  additionally a constant $\Ihat{c}^0$ (the \emph{pointer} of $c$), for
  each term definition $c \defeq t : \tau$ of $\Env$, where the
  \emph{erasure arity} $a_c$ is determined from $\erase{t}$ in
  \cref{def:compile};
\item a binary function symbol $\fapp$ (written infix and associating to
  the left), and a constant $\prfc$ (the image of dead code $\bx$);
\item countably many \emph{auxiliary} function symbols, produced by the
  compilation procedure below.  As in \cite[\S 4]{Czajka2018} we assume
  injective naming functions $\Lambda(\cdot)$, $\mathrm{X}(\cdot)$,
  $\Phi(\cdot)$ from $\lbox$ terms (up to renaming of free variables) to
  auxiliary symbols, so that the symbol lifted for a given term is unique
  and terms differing only in variable names share their symbol
  (hash-consing).
\end{itemize}
Source term variables double as first-order variables.
\end{definition}

There is \emph{no} binary inhabitation predicate in $\SigFOL$: the only
typing-flavoured atoms are the per-datatype unary guards
$\Ihat{I}(\cdot)$, and they occur only at atomic-type leaves of the
translation.

\subsection{Compilation of erased programs}
\label{sec:translation-compile}

The compilation turns a closed $\lbox$ program into defining equations
with constructor patterns.  It is presented as a recursive procedure that
\emph{emits} first-order axioms and, for every symbol $f$ it introduces,
records a \emph{witness}: a $\lbox$ term $e_f$ with free variables among
the closure parameters of $f$, used by the model construction of
\cref{sec:semantics} and nowhere else.

Compilation carries a finite map $\theta$ assigning to certain $\lbox$
variables (fix-bound recursion variables) a \emph{template}
$(f, \tup{w}, m)$: an $(|\tup{w}| + m)$-ary symbol $f$, a list of
first-order terms $\tup{w}$, and an argument count $m$.

\begin{definition}[Term translation $\compf_\theta$]
\label{def:compC}
For a $\lbox$ term $e$ whose free variables are first-order variables or
$\theta$-bound, the first-order term $\compf_\theta(e)$ is defined by
recursion on $e$:
\begin{itemize}[leftmargin=2em]
\item $\compf_\theta(x) = x$ for an ordinary variable;
  $\compf_\theta(\bx) = \prfc$;
  $\compf_\theta(\wcns{c}(\tup{e})) =
    \Ihat{\wcns{c}}(\compf_\theta(\tup{e}))$.
\item For an application spine $e = h\,e_1 \cdots e_j$ ($j \geq 0$, $h$
  not an application, $e$ itself not of the previous forms):
  \begin{itemize}
  \item if $h = g$ with $\theta(g) = (f, \tup{w}, m)$: by the guard
    condition $j \geq m$, and
    $\compf_\theta(e) = f(\tup{w}, \compf_\theta(e_1), \dots,
    \compf_\theta(e_m)) \fapp \compf_\theta(e_{m+1}) \fapp \cdots \fapp
    \compf_\theta(e_j)$;
  \item if $h = c$ is a term-definition constant: if $j \geq a_c$ then
    $\compf_\theta(e) = \Ihat{c}(\compf_\theta(e_1), \dots,
    \compf_\theta(e_{a_c})) \fapp \cdots \fapp \compf_\theta(e_j)$, and
    otherwise $\compf_\theta(e) = \Ihat{c}^0 \fapp \compf_\theta(e_1)
    \fapp \cdots \fapp \compf_\theta(e_j)$;
  \item if $h = x$ is a variable: $\compf_\theta(e) = x \fapp
    \compf_\theta(e_1) \fapp \cdots \fapp \compf_\theta(e_j)$;
  \item if $h$ is a $\lambda$-abstraction, case or fix (a \emph{complex
    head}, necessarily with $j \geq 0$): let $u$ be its \emph{value
    translation} defined below; then $\compf_\theta(e) = u \fapp
    \compf_\theta(e_1) \fapp \cdots \fapp \compf_\theta(e_j)$.
  \end{itemize}
\item Value translations of complex heads $h$ with $\tup{w} =
  \fv{h}$ (ordered):
  \begin{itemize}
  \item $h = \lambda x.e_0$: let $f = \Lambda(h)$, an $|\tup{w}|$-ary
    symbol with witness $e_f = h$; run
    $\mathrm{B}_\theta(f(\tup{w});\, h)$; the value is $f(\tup{w})$.
  \item $h = \mathsf{case}(s;\, \tup{br})$: let $g = \mathrm{X}(\lambda
    z.\mathsf{case}(z; \tup{br}))$ for fresh $z$, of arity $1 +
    |\tup{w}'|$ with $\tup{w}' = \fv{\tup{br}}$, with witness $e_g =
    \lambda z.\mathsf{case}(z; \tup{br})$; run
    $\mathrm{B}_\theta(g(z, \tup{w}');\, \mathsf{case}(z; \tup{br}))$; the
    value is $g(\compf_\theta(s), \tup{w}')$.
  \item $h = \mathsf{fix}\,g_0.\,\lambda z_1 \dots \lambda z_m.\,e_0$
    (maximal prefix): let $f = \Phi(h)$ of arity $|\tup{w}| + m$ with
    witness $e_f = h$, and $f^0 = \Phi^0(h)$ of arity $|\tup{w}|$ with
    witness $e_{f^0} = h$; emit the pointer axiom $\forall \tup{w}\,
    \tup{z}.\ f^0(\tup{w}) \fapp z_1 \fapp \cdots \fapp z_m \feq
    f(\tup{w}, \tup{z})$; run $\mathrm{B}_{\theta'}(f(\tup{w},
    \tup{z});\, e_0)$ where $\theta' = \theta[g_0 \mapsto (f, \tup{w},
    m)]$; the value is $f^0(\tup{w})$.
  \end{itemize}
\end{itemize}
\end{definition}

\begin{definition}[Body compilation $\mathrm{B}$]
\label{def:compile}
$\mathrm{B}_\theta(L;\, e)$ takes a first-order term $L$ (the
\emph{left-hand side}, linear in its variables, whose variables include
$\fv{e}$ minus the $\theta$-bound ones) and a $\lbox$ term $e$, and emits
axioms:
\begin{enumerate}[leftmargin=2em]
\item If $e = \lambda x.e'$: run $\mathrm{B}_\theta(L \fapp x;\, e')$.
\item If $e = \mathsf{case}(x;\, \{\wcns{c}_i(\tup{z}_i) \Rightarrow
  b_i\}_i)$ with $x$ a variable of $L$: for each $i$, with $\tup{z}_i$
  fresh, run
  $\mathrm{B}_\theta\bigl(
    \subst{L}{\Ihat{\wcns{c}}_i(\tup{z}_i)}{x};\;
    \subst{b_i}{\wcns{c}_i(\tup{z}_i)}{x}\bigr)$
  --- the left-hand side is refined by the (hatted) constructor pattern,
  and the remaining occurrences of the scrutinee variable in the branch
  are replaced by the (source) constructor pattern.
\item If $e = \mathsf{case}(s;\, \tup{br})$ with $s$ not a variable of
  $L$: let $g = \mathrm{X}(\lambda z.\mathsf{case}(z;\tup{br}))$ with
  closure $\tup{w}' = \fv{\tup{br}}$ and witness as in
  \cref{def:compC}; emit
  \[
    \forall.\ L \feq g(\compf_\theta(s), \tup{w}')
  \]
  (universal closure over the variables of $L$), and run
  $\mathrm{B}_\theta(g(z, \tup{w}');\, \mathsf{case}(z; \tup{br}))$ for
  fresh $z$.
\item If $e = \mathsf{fix}\,g_0.e'$: compute the value translation $u =
  f^0(\tup{w})$ of $e$ as in \cref{def:compC} (which also compiles $e$);
  emit $\forall.\ L \feq u$.
\item Otherwise ($e$ not a $\lambda$, case or fix at the root): emit
  $\forall.\ L \feq \compf_\theta(e)$.
\end{enumerate}
For a term definition $c \defeq t : \tau$ of $\Env$, compilation is
initialised as follows on $e = \erase{t}$ (a closed $\lbox$ term):
\begin{itemize}[leftmargin=2em]
\item if $e = \mathsf{fix}\,g_0.\lambda z_1 \dots \lambda z_m.e_0$
  (maximal prefix): set $a_c \defeq m$ and run
  $\mathrm{B}_{[g_0 \mapsto (\Ihat{c},\, (),\, m)]}(\Ihat{c}(\tup{z});\,
  e_0)$;
\item otherwise, with $e = \lambda z_1 \dots \lambda z_m.e_0$ (maximal
  prefix, possibly $m = 0$): set $a_c \defeq m$ and run
  $\mathrm{B}_{\emptyset}(\Ihat{c}(\tup{z});\, e_0)$.
\end{itemize}
The witness of $\Ihat{c}$ and of $\Ihat{c}^0$ is the $\lbox$ constant $c$;
if $a_c > 0$, the pointer axiom
$\forall \tup{z}.\ \Ihat{c}^0 \fapp z_1 \fapp \cdots \fapp z_{a_c} \feq
\Ihat{c}(\tup{z})$ is emitted.
\end{definition}

\begin{lemma}[Compilation terminates]
\label{lem:compile-term}
The procedure of \cref{def:compC,def:compile} terminates on every input,
and (with the naming functions of \cref{def:signature}) emits a finite set
of axioms for a given environment and goal.
\end{lemma}

\begin{proof}
Define the size $|e|$ of a $\lbox$ term as its number of syntax nodes.
Every recursive call of $\mathrm{B}$ or $\compf$ is on a term of strictly
smaller size, with one exception: rule (3) of $\mathrm{B}$ passes from
$\mathsf{case}(s; \tup{br})$ to $\mathsf{case}(z; \tup{br})$.  Since $s$
is not a variable in rule (3), $|s| > 1 = |z|$, so the size still strictly
decreases.  In rule (2) the branch $b_i$ has its scrutinee variable
replaced by a constructor pattern of fresh \emph{variables}; the
substituted occurrences of $x$ in $b_i$ are variable occurrences replaced
by terms of size $1 + r$ --- to make the measure genuinely decrease we
count the size of the \emph{pair} (number of case/λ/fix nodes of $e$, node
count), ordered lexicographically: rule (2) consumes the outer case node
and passes to a proper subterm $b_i$ of $e$ up to a variable-for-pattern
replacement that does not create case, $\lambda$ or fix nodes; rules (1),
(3), (4), (5) and all $\compf$ recursions also strictly decrease the pair.
Finiteness follows since each emitted axiom is attached to a symbol, each
symbol to a $\lbox$ term occurring in the (finite) compilation tree, and
hash-consing merges repeated liftings.
\end{proof}

\begin{remark}[Shape of the equations]
\label{rem:eq-shape}
All axioms emitted by compilation are universally quantified equations
$L \feq R$ whose left-hand sides are linear terms of the form
$f(\tup{\pi}) \fapp x_1 \fapp \cdots \fapp x_k$ with $\tup{\pi}$
constructor patterns and variables.  Grouped by their head symbol, the
left-hand sides arising from one case tree are pairwise non-overlapping
(they refine disjoint constructor patterns), so the emitted equations form
an \emph{orthogonal} first-order rewrite system when oriented left to
right \cite{BaaderNipkow1998}.  This property is not used in the soundness
proof, but it is the operational point of the exercise: the equations are
orientable rewrite rules, usable by superposition provers for demodulation
and by SMT solvers as triggers (\cref{sec:casesplit-cnf}).
\end{remark}

\begin{example}[Basic definitions]
\label{ex:compile}
Let $\nat$ have constructors $\csO : ()$ and $\csS : (\nat)$, and $\bool$
have $\cstrue, \csfalse : ()$.  Writing the hatted symbols without hats
for readability:
\begin{enumerate}[leftmargin=2em]
\item $\mathsf{mypred} \defeq \lambda x{:}\nat.\,
  \tcase{x}{\_.\nat}{\csO \Rightarrow \csO \mid \csS(y) \Rightarrow y}$
  compiles to
  \[
    \mathsf{mypred}(\csO) \feq \csO, \qquad
    \forall y.\ \mathsf{mypred}(\csS(y)) \feq y.
  \]
\item $\mathsf{add} \defeq \tfix\,f{:}\Pi n\,m{:}\nat.\nat.\;
  \lambda n\,m.\,\tcase{n}{\_.\nat}{\csO \Rightarrow m \mid \csS(y)
  \Rightarrow \csS(f\,y\,m)}$ compiles to
  \[
    \forall m.\ \mathsf{add}(\csO, m) \feq m, \qquad
    \forall y\,m.\ \mathsf{add}(\csS(y), m) \feq
      \csS(\mathsf{add}(y, m)).
  \]
\item Nested matches flatten into deeper patterns:
  $g\,x \defeq \tcase{x}{}{\csO \Rightarrow \csO \mid \csS(y) \Rightarrow
  \tcase{y}{}{\csO \Rightarrow \csS(\csO) \mid \csS(z) \Rightarrow z}}$
  compiles to $g(\csO) \feq \csO$, $g(\csS(\csO)) \feq \csS(\csO)$ and
  $\forall z.\ g(\csS(\csS(z))) \feq z$.
\item A compound scrutinee is flattened through an auxiliary symbol:
  $k\,b \defeq \tcase{\mathsf{negb}\,b}{}{\cstrue \Rightarrow \csO \mid
  \csfalse \Rightarrow \csS(\csO)}$ compiles to
  \[
    \forall b.\ k(b) \feq g'(\mathsf{negb}(b)), \qquad
    g'(\cstrue) \feq \csO, \qquad
    g'(\csfalse) \feq \csS(\csO),
  \]
  where $g'$ is the auxiliary symbol lifted for the case expression.
\end{enumerate}
No axiom in this list carries a guard: compare
\cref{sec:casesplit}.
\end{example}

\subsection{Specification extraction}
\label{sec:translation-spec}

\begin{definition}[Guard formulas and proposition translation]
\label{def:guard-F}
For every shallow set type $\sigma$ and first-order term $u$, the
\emph{guard formula} $\guard{\sigma}{u}$, and for every shallow
proposition $\varphi$, the first-order formula $\F{\varphi}$, are defined
by mutual structural recursion:
\begin{align*}
  \guard{I}{u} &= \Ihat{I}(u)\\
  \guard{\Pi x{:}\sigma.\tau}{u} &=
    \forall x.\ \guard{\sigma}{x} \to \guard{\tau}{u \fapp x}\\
  \guard{\varphi \to \tau}{u} &= \F{\varphi} \to \guard{\tau}{u}\\
  \guard{\refty{x{:}\sigma}{\varphi}}{u} &=
    \guard{\sigma}{u} \wedge \subst{\F{\varphi}}{u}{x}
  \\[1ex]
  \F{\top} = \top, \quad \F{\bot} &= \bot, \quad
  \F{\varphi \circledast \psi} = \F{\varphi} \circledast \F{\psi}
  \quad (\circledast \in \{\wedge, \vee, \to\})\\
  \F{\forall x{:}\sigma.\varphi} &=
    \forall x.\ \guard{\sigma}{x} \to \F{\varphi}\\
  \F{\exists x{:}\sigma.\varphi} &=
    \exists x.\ \guard{\sigma}{x} \wedge \F{\varphi}\\
  \F{a =_{\sigma} b} &= \compf(\erase{a}) \feq \compf(\erase{b})\\
  \F{J(\tup{a})} &= \Ihat{J}(\compf(\erase{a_1}), \dots,
    \compf(\erase{a_k}))
\end{align*}
where in the clauses for $\Pi$, $\forall$, $\exists$ and refinement the
bound source variable is reused as the bound first-order variable, and
$\compf = \compf_\emptyset$ is the term translation of \cref{def:compC}
(whose complex-head clauses may emit further axioms).  Note that
$\guardf$ and $\Ff$ are total on $\Sfrag$ and structurally recursive; no
normalisation of the source is involved.
\end{definition}

Two clauses implement the shallowness requirements discussed in
\cref{sec:intro,sec:shallow}.  The refinement clause expands the
propositional payload \emph{inline at each occurrence} --- positively
under $\exists$ and in conclusions, negatively under $\forall$ and in
premises --- rather than emitting an opaque atom plus an unfolding axiom.
The premise clause $\guard{\varphi \to \tau}{u} = \F{\varphi} \to
\guard{\tau}{u}$ does \emph{not} apply $u$ to anything: erasure has
dropped the proof argument, so the specification tracks the pruned arity
of the extracted program.

\begin{definition}[Specification of a constant]
\label{def:spec}
For a term definition $c \defeq t : \tau$ with $\tau \in \Sfrag$, the
\emph{specification axiom} $\specf(c)$ is $S(u_0; \tau)$ where $u_0 =
\Ihat{c}$ if $a_c = 0$ and otherwise $u_0 = \Ihat{c}(\,)$ is the empty
application accumulator, and
\begin{align*}
  S(u; \Pi x{:}\sigma.\tau') &= \forall x.\ \guard{\sigma}{x} \to
    S(u \cdot x;\, \tau')\\
  S(u; \varphi \to \tau') &= \F{\varphi} \to S(u; \tau')\\
  S(u; \tau') &= \guard{\tau'}{\overline{u}} \quad
    \text{if $\tau'$ is a datatype or refinement type,}
\end{align*}
where $u \cdot x$ appends $x$ to the accumulator if fewer than $a_c$
arguments have been accumulated and otherwise applies with $\fapp$, and
$\overline{u}$ closes the accumulator: an under-applied accumulator
$\Ihat{c}(x_1, \dots, x_j)$ with $j < a_c$ is rendered as the pointer
chain $\Ihat{c}^0 \fapp x_1 \fapp \cdots \fapp x_j$.
\end{definition}

\begin{example}[The running example $h$]
\label{ex:h}
Consider
\begin{gather*}
  h \defeq \lambda x\,y\,z{:}\nat.\ \lambda p{:}(x =_\nat y \wedge y
  =_\nat z).\ \tsplit{p}{p_1\,p_2.\ \trefine{z}{\,\mathsf{trans}(p_1,
  p_2)}}\\
  h \;:\; \Pi x\,y\,z{:}\nat.\ (x = y \wedge y = z) \to
  \refty{u{:}\nat}{x =_\nat u}
\end{gather*}
where $\mathsf{trans}(p_1,p_2)$ abbreviates the evident proof of
$x =_\nat z$ by $\mathsf{rew}$.  Erasure gives
\[
  \erase{h} = \lambda x\,y\,z.\,z,
\]
(the proof abstraction disappears, $\tsplit{}{}$ collapses to its branch,
$\trefine{}{}$ unboxes), so $a_h = 3$ and compilation emits the single
equation
\[
  \forall x\,y\,z.\ \Ihat{h}(x, y, z) \feq z.
\]
Specification extraction on the type gives
\begin{multline*}
  \specf(h) =
  \forall x.\,\nat(x) \to \forall y.\,\nat(y) \to \forall z.\,\nat(z)
  \to{}\\
  \Bigl( (x \feq y \wedge y \feq z) \to
    \bigl( \nat(\Ihat{h}(x,y,z)) \wedge x \feq \Ihat{h}(x,y,z) \bigr)
  \Bigr).
\end{multline*}
Both axioms requested by the design are produced, and the symbols
$\mathsf{sig}$, $\mathsf{exist}$, $\mathsf{proj1\_sig}$ and any
inhabitation atom are absent from the output.  Note the interplay of the
components: the equation is blind to the propositional content ($h$ is
extractionally a projection), while the specification carries it; for an
underspecified function the situation is reversed.
\end{example}

\subsection{The theory \texorpdfstring{$\ThE$}{Delta(E,G)}}
\label{sec:translation-theory}

\begin{definition}[Emitted theory]
\label{def:theory}
The first-order theory $\ThE$ consists of the following axioms.
\begin{enumerate}[leftmargin=2em]
\item For each datatype $I$ with constructors $\wcns{c}_i :
  (\sigma_{i1}, \dots, \sigma_{ir_i})$:
  \begin{itemize}
  \item (\emph{constructor typing})\ \
    $\forall \tup{y}.\ \bigwedge_j \guard{\sigma_{ij}}{y_j} \to
     \Ihat{I}(\Ihat{\wcns{c}}_i(\tup{y}))$ \ for each $i$;
  \item (\emph{inversion})\ \
    $\forall w.\ \Ihat{I}(w) \to \bigvee_i \exists \tup{y}.\
     \bigwedge_j \guard{\sigma_{ij}}{y_j} \wedge
     w \feq \Ihat{\wcns{c}}_i(\tup{y})$;
  \item (\emph{injectivity})\ \
    $\forall \tup{y}\,\tup{y}'.\ \Ihat{\wcns{c}}_i(\tup{y}) \feq
     \Ihat{\wcns{c}}_i(\tup{y}') \to \bigwedge_j y_j \feq y'_j$ \ for
     each $i$ with $r_i > 0$;
  \item (\emph{discrimination})\ \
    $\forall \tup{y}\,\tup{y}'.\ \neg\bigl(\Ihat{\wcns{c}}_i(\tup{y})
     \feq \Ihat{\wcns{c}}_j(\tup{y}')\bigr)$ \ for all $i \neq j$.
  \end{itemize}
\item For each inductive predicate $J$ with clauses $\wcns{k}_j : \forall
  \tup{x}{:}\tup{\tau}_j.\ \varphi_{j1} \to \dots \to \varphi_{j\ell_j}
  \to J(\tup{u}_j)$:
  \begin{itemize}
  \item (\emph{clause})\ \
    $\forall \tup{x}.\ \bigwedge_i \guard{\tau_{ji}}{x_i} \to
     \F{\varphi_{j1}} \to \dots \to \F{\varphi_{j\ell_j}} \to
     \Ihat{J}(\compf(\erase{\tup{u}_j}))$;
  \item (\emph{case analysis})\ \
    $\forall \tup{z}.\ \Ihat{J}(\tup{z}) \to \bigvee_j \exists
     \tup{x}.\ \bigwedge_i \guard{\tau_{ji}}{x_i} \wedge \bigwedge_i
     \F{\varphi_{ji}} \wedge \bigwedge_l z_l \feq
     \compf(\erase{u_{jl}})$.
  \end{itemize}
\item For each term definition $c \defeq t : \tau$: the equations, pointer
  axioms and auxiliary axioms emitted by compiling $\erase{t}$
  (\cref{def:compile}), and the specification axiom $\specf(c)$ if $\tau
  \in \Sfrag$.
\item For each lemma $c \defeq \pi : \varphi$ with $\varphi \in \Sfrag$:
  the \emph{premise axiom} $\F{\varphi}$.
\item The auxiliary axioms emitted by the computation of $\F{G}$ (via the
  complex-head clauses of $\compf$ on terms embedded in $G$).
\end{enumerate}
The \emph{translation of the problem} $(\Env, G)$ is the pair $(\ThE,
\F{G})$.
\end{definition}

\begin{remark}
\label{rem:theory-notes}
(i) The case-analysis axiom for an inductive predicate is the
formula-level image of case analysis on a derivation --- inversion in the
proof-assistant sense; full induction is a schema over motives and is not
emitted (an ATP would rarely be able to instantiate it; when induction is
needed it is the reconstruction backend's job).  (ii) Premise selection
operates before translation, so in practice $\ThE$ is generated from the
selected slice of the environment; soundness is monotone in the axiom set,
so \cref{thm:soundness} covers every subset of $\ThE$.  (iii) The
per-datatype inversion axiom is where exhaustiveness of constructors
lives; the defining equations of \cref{def:compile} make no such claim.
This division of labour is the subject of \cref{sec:casesplit}.
\end{remark}
